\begin{table}[t]
\centering
\small
\begin{tabular}{p{2.2cm} p{3.1cm} p{1.4cm}}
\toprule
 \textbf{Reference} & \textbf{Model output} & \textbf{MSA feature}\\
\midrule
  {{vocalized_transcript 7896_1}}
    \newline \scriptsize {{subdialect 7896_1}}
  & Catt: {{diacritization/catt 7896_1 bold=كَيْفَ,يُعْرَفُ}}
    \newline Gemini: {{diacritization/gemini-3.5-flash 7896_1}}
    & case endings \\
  {{vocalized_transcript LEB_120}}
    \newline \scriptsize {{subdialect LEB_120}}
  & Mishkal: {{diacritization/mishkal LEB_120 bold=كُلّهَا}}
    \newline Gemini: {{diacritization/gemini-3.5-flash LEB_120}}
    & MSA vowels \\
  {{vocalized_transcript LEB_0 :3}}
    \newline \scriptsize {{subdialect LEB_0}}
  & Mishkal: {{diacritization/mishkal LEB_0 bold=نُحْنَا,بِمَوْسِمٍ :3}}
    \newline Levanti: {{diacritization/levanti LEB_0 :3}}
    & short vowels \\
    \bottomrule
\end{tabular}
\caption{Diacritization examples: an MSA-oriented dedicated system (Catt,
Mishkal) vs.\ a Levantine-aware system (Gemini, Levanti) on the same utterance.
\textbf{Bold} marks words that still differ from the reference after the
normalization of \S\ref{sec:app_normalization}, so sukun and the MSA
case-\textit{tanwin} are not marked; the Modern Standard Arabic feature that
the dedicated diacritizer restores but Levantine speech does not is named at
right, and the Levantine-aware system produces the dialectal form.}
\label{tab:diac-qualitative}
\end{table}
